%% Supplement 1 to the imputation_article main manuscript.
%% Topic: alternatives and supplements to ComBat for integrating
%%        microarray studies whose biological groups (here pregnancy
%%        trimesters) are not all represented in every dataset.
%% Compile: pdflatex supplement1.tex (twice, for refs)

\documentclass[11pt,a4paper]{article}

\usepackage[margin=1in]{geometry}
\usepackage{amsmath}
\usepackage{booktabs}
\usepackage{array}
\usepackage{longtable}
\usepackage{url}
\usepackage{xcolor}
\usepackage{hyperref}
\hypersetup{colorlinks=true,linkcolor=blue,citecolor=blue,urlcolor=blue}

\newcommand{\doi}[1]{\href{https://doi.org/#1}{doi:#1}}

\title{Supplement 1: Alternatives and supplements to ComBat for
  integrating microarray studies with non-overlapping
  biological groups}
\author{}
\date{}

\begin{document}
\maketitle

\section*{Motivation}

ComBat~\cite{johnson2007} is the default empirical-Bayes method for
removing batch effects from merged microarray data and is the engine
used in the main pipeline. Its location-and-scale model assumes that
each batch contributes additive and multiplicative perturbations that
are \emph{independent} of the biological covariates of interest. When
the covariate of interest is balanced across batches, this assumption
is harmless: the per-batch mean of the biological signal is the same,
so subtracting per-batch means does not erase between-group
contrasts.

The placental dataset collection used in this work violates that
assumption in a hard way. Several Affymetrix and Illumina studies
contain only a single trimester (e.g.\ first trimester only, or term
only); other studies span two trimesters but never all three.
There is no batch in which all three trimesters are observed jointly.
Under this design, ``batch'' is partially aliased with ``trimester'':
the batch-mean estimate that ComBat removes contains real biological
signal, and the model has no way to separate the two without
additional information.

Nygaard, R{\o}dland, and Hovig formalised the problem and showed
that running ComBat (or its parametric two-way ANOVA equivalent) on
unbalanced designs and then handing the corrected matrix to a naive
downstream test exaggerates between-group differences and inflates
false positives, sometimes by an order of magnitude~\cite{nygaard2016}.
Their concrete recommendation is to skip the pre-correction step and
instead include batch as a covariate in the differential-expression
model directly, so that the residual degrees of freedom are charged
correctly.

Towfic, Kusko, and Zeskind replied with a re-analysis of one of the
datasets used by Nygaard (GSE61901), arguing that the inflation
Nygaard reported was driven not by ComBat itself but by how
\emph{technical replicates} were handled in the analysis pipeline.
When technical replicates were retained and modelled with limma's
\texttt{duplicateCorrelation} or with mixed-effects regression
rather than averaged away, several batch-handling strategies
(including ComBat followed by limma) converged on a similar
$\sim\!1500$-DEG result, while the averaging strategy collapsed it
to 11 DEGs~\cite{towfic2017}. The two papers therefore agree on the
core diagnostic --- when the design is unbalanced the choice of
batch-handling pipeline materially changes the differential
expression list and must be documented --- but they disagree about
which step in the pipeline is to blame: Nygaard fingers the
ComBat-then-test sequence, Towfic et al.\ fingers the averaging of
technical replicates.

This supplement reviews the methods that are most commonly used as
\emph{alternatives} or \emph{supplements} to ComBat for this regime,
groups them by the assumption they substitute for the
``balanced design'' assumption, and gives practical recommendations
for the placenta integration use case. Downloaded references are
listed in
\texttt{references/integration\_methods\_references/}.

%% =====================================================================
\section{Why ComBat is fragile on unbalanced trimester data}

ComBat fits, for every gene $g$ and every sample $i$ in batch $b$:
\begin{equation}
  Y_{gi} \;=\; \alpha_g \;+\; X_i\,\beta_g \;+\; \gamma_{gb} \;+\;
  \delta_{gb}\,\varepsilon_{gi},
  \label{eq:combat}
\end{equation}
where $X_i$ encodes the biological covariates that should be
preserved, $\gamma_{gb}$ is the additive batch shift, and
$\delta_{gb}$ is the multiplicative batch scale. Empirical-Bayes
shrinkage of $\gamma_{gb}$ and $\delta_{gb}$ across genes makes the
fit stable even with few samples per batch.

The crucial step is that ComBat then \emph{subtracts} the estimated
$\hat{\gamma}_{gb}$ and divides by $\hat{\delta}_{gb}$ to produce a
``corrected'' $Y^{*}_{gi}$ that downstream methods treat as if it
came from a single batch. The estimate $\hat{\gamma}_{gb}$ is the
within-batch residual mean after removing $X_i\,\beta_g$. If batch
$b$ contains only first-trimester samples, then ``the within-batch
mean after accounting for biology'' is meaningless: the design matrix
cannot identify $\beta_g$ from $\gamma_{gb}$ inside that batch alone,
and the cross-batch pooling of $\hat{\gamma}_{gb}$ ends up importing
trimester structure into the batch shift. The result is that the
trimester effect is partially absorbed into the batch correction and
then subtracted off.

Two observations follow:
\begin{enumerate}
  \item If the goal is ``one merged matrix that PCA, clustering, and
        visualisation can use,'' the bias from unbalanced ComBat is
        irreducible without strong external assumptions.
  \item If the goal is differential expression, it is almost always
        safer to keep batch in the design matrix of the
        downstream linear model and \emph{not} pre-subtract a batch
        shift~\cite{nygaard2016,towfic2017}.
\end{enumerate}

The methods reviewed below trade the ``balanced design'' assumption
for one of: (a) the existence of negative-control genes, (b) the
existence of a reference batch or replicate samples, (c) the
existence of mutual-nearest-neighbour structure between batches, or
(d) the willingness to combine effect sizes per study rather than
samples.

%% =====================================================================
\section{Alternatives and supplements}

\subsection{Surrogate Variable Analysis (SVA)}

SVA~\cite{leek2007} sidesteps the problem of \emph{naming} batches.
Instead of asking the user which samples share a batch, it estimates
latent factors (``surrogate variables'') from the residuals after
removing the biological model, and adds those factors as covariates
in the downstream test. Because the latent factors are estimated
\emph{conditional on the protected biological covariate}, SVA is in
principle robust to designs where the explicit batch label is
confounded with biology.

In practice SVA is very useful when the technical structure is
multi-dimensional (lab effect $+$ scan date $+$ ozone). Leek and
Storey themselves note that SVA's identifiability rests on the
assumption that a substantial subset of genes is unaffected by the
biological covariate; if that assumption is violated, the surrogate
variables can become correlated with biology and remove signal
along with noise. Although the original paper does not work out
this failure mode for the specific case of a batch that is fully
aliased with biology, the consequence is direct: with placenta
data where one batch contains a single trimester, the leading
residual direction is essentially ``this study,'' and any
data-driven surrogate-variable estimator is at risk of capturing
trimester structure inside its surrogate factor. We therefore
recommend that on this collection SVA be used \emph{instead of}
ComBat (never on top of it) and that the recovered surrogate
variables always be plotted against the trimester labels before
being trusted.

\subsection{Remove Unwanted Variation (RUV)}

The RUV family~\cite{gagnonbartsch2012,risso2014} replaces the
``balanced design'' assumption with ``the user can name a set of
genes that are not differentially expressed with respect to the
biological covariate of interest.'' Variation in those negative
control genes is by construction technical, so a factor analysis
restricted to them yields nuisance factors that can be added to the
limma design.

RUV-2 (RUVg) uses negative-control genes; RUV-4 and RUVs allow
control \emph{samples} (e.g.\ technical replicates) instead.
Importantly RUV does \emph{not} require that controls be present in
every batch: the factor model is estimated globally. This makes RUV
the most direct ``ComBat replacement'' for unbalanced placental
data, provided a sensible negative-control list can be assembled.
The validity of that list is the load-bearing assumption of the
whole approach --- Gagnon-Bartsch and Speed are explicit that
negative-control genes must be ones whose expression is known
\emph{a priori} not to depend on the biological variable of
interest, and the method gives no diagnostic if that assumption is
broken. For trimester comparisons a reasonable starting point is
housekeeping genes from the Eisenberg/Levanon list intersected with
genes that have low variance across the trimester axis in a single
balanced study (e.g.\ a study that contains both endpoints of the
contrast); any list assembled this way must then be validated
against an independent balanced subset before being trusted.

A pragmatic recipe:
\begin{enumerate}
  \item Pick negative-control genes (housekeeping, or empirically
        invariant in a balanced reference study).
  \item Run \texttt{RUVg(set, controls, k)} with $k$ chosen by the
        elbow of the singular values, capped at the number of
        batches minus one.
  \item Add the \texttt{W} factors as covariates to the limma
        design alongside the trimester contrast.
  \item Plot \texttt{W} against trimester to verify that the
        surrogate factors are \emph{not} just trimester re-labelled.
\end{enumerate}

\subsection{Reference-batch ComBat (BESC, ComBat-ref)}

A second family of methods keeps the ComBat additive/multiplicative
model but eliminates the within-batch identifiability problem by
nominating one batch as a \emph{reference}.

Walker et al.~\cite{walker2008} were the earliest concrete proposal:
include identical pooled-RNA reference samples, derived from the
same tissue as the biological samples, in every batch. Their
correction is essentially ComBat with the reference samples used
to anchor each batch's location and scale, allowing biological
groups to be unbalanced across batches as long as the reference
pool is shared.

Varma's BESC~\cite{varma2020} addresses the same problem
independently. Rather than requiring physical reference samples in
every batch, BESC precomputes a small set of ``batch effect
signatures'' (vectors of length \(p\), the number of probes) on a
large public reference cohort, and at correction time projects
each new sample onto those signatures and subtracts them. Because
the signatures are estimated once and then applied per sample,
BESC is conservative (it only removes variation in directions that
were technical in the reference cohort) but cannot remove batch
effects whose direction was not represented there.

ComBat-ref~\cite{zhang2025} is the most recent entry. It fits the
negative-binomial ComBat-seq model~\cite{zhang2020} on count data, but instead of
treating all batches symmetrically it estimates a per-batch
dispersion \(\hat{\lambda}_b\) and \emph{automatically selects the
batch with the smallest dispersion as the reference}. Counts in
the reference batch are then preserved unchanged, and every other
batch is adjusted towards it. This automation removes the need for
the user to nominate a reference batch by hand. ComBat-ref is
implemented as a small set of R scripts at
\url{https://github.com/xiaoyu12/Combat-ref} (a copy of the
sources is included in
\texttt{references/integration\_methods\_references/}).

These methods are well suited to the placenta problem if a
reference batch can be chosen --- whether by hand or, as in
ComBat-ref, automatically --- that contains samples from a broad
range of trimesters and has stable internal variance. In practice
the closest reference candidate in the current collection is the
dataset that spans the broadest sample of trimesters. We recommend
testing reference-batch ComBat with that dataset as the anchor and
comparing the resulting differential expression list against the
imputation pipeline output. ComBat-ref itself targets RNA-seq
counts; for the existing log-transformed microarray matrices the
relevant analogues are Walker-style anchored ComBat or BESC.

\subsection{Mutual Nearest Neighbours (MNN, fastMNN, scMerge, Harmony)}

MNN~\cite{haghverdi2018} and Harmony~\cite{korsunsky2019} were
developed for single-cell RNA-seq integration but they are agnostic
about the data type and have been used on bulk
microarrays~\cite{varma2020}. The core idea is that batches are
\emph{aligned} rather than \emph{normalised}: pairs of samples
across batches that are mutual nearest neighbours in expression
space are taken as evidence that those samples represent the same
biological state, and a translation vector is computed to bring
them into register.

MNN-style alignment is attractive for the placenta integration
problem because it requires \emph{some} biological overlap between
batches but not the full factorial: as long as one or two samples
per pair of studies represent comparable trimesters, the alignment
can stitch the studies together. The downside is that MNN was
designed for thousands of cells, not for tens of bulk samples, and
the variance estimates that drive the neighbour search are noisy
on bulk data. We therefore recommend MNN only as a sanity check on
top of ComBat or RUV, not as the sole correction.

\subsection{Direct integration of cross-platform intensity data}

Turnbull et al.~\cite{turnbull2012} demonstrated that intensity
levels from Affymetrix and Illumina single-channel arrays can be
merged at the probe level after within-platform normalisation, and
that the merged data have higher statistical power than any single
platform alone. They benchmarked several normalisation /
batch-correction combinations and found that ComBat and
cross-platform normalisation (XPN) outperformed mean-centring and
distance-weighted discrimination (DWD); their workflow allows
ComBat as an optional step when batch effects remain large after
within-platform normalisation. The takeaway for the placenta
collection is therefore not ``skip ComBat'' but rather ``do not
assume that a single normalisation step is enough --- benchmark
several combinations and pick the one that minimises residual
batch structure in PCA.''

Almeida-de-Macedo et al.~\cite{macedo2013} studied a related
question: when correlation coefficients between gene pairs are
estimated from heterogeneous microarray data, is it better to pool
the data first (``melting pot'') or to estimate the correlation in
each study and combine afterwards (``mosaic'')? Their primary
finding is about bias in correlation estimates rather than about
DE detection: pooled estimates can be substantially biased when
group sizes are heterogeneous across studies, whereas the mosaic
combination of per-study estimates is less biased at the cost of
statistical power. The paper does not directly address
batch-confounded differential expression, but its conclusion ---
that combination of per-study summaries is safer than direct
pooling when studies are heterogeneous --- carries over and
motivates the effect-size meta-analysis approach in the next
section.

\subsection{Effect-size meta-analysis (the mosaic approach)}

Meta-analytic combination of per-study effect sizes does not need
any cross-batch normalisation at all: each study contributes a
$\log\!\mathit{FC}$ and a standard error, and a random-effects
model combines them. Crucially, the per-study contrast can only be
computed in studies that contain \emph{both} of its endpoints, so
the unbalanced studies that motivate this supplement contribute to
some contrasts but not others.

A related but distinct problem is studied by Shi et
al.~\cite{shi2011}: their ``meta-analysis with missing
replicates'' framework targets \emph{within-study} missing data
(genes that are absent on some platforms or replicates that
failed) and combines per-study summaries with imputation of the
missing entries. The methodology is not a one-to-one match for
the placenta case, where the missingness is at the level of
\emph{whole biological groups} per study rather than of individual
replicates, but the same effect-size combination machinery
(\texttt{metafor}-style random-effects models) applies. For the
placenta collection, an effect-size meta-analysis should be
reported \emph{alongside} the ComBat / imputation pipeline as a
sensitivity analysis: any gene that is differentially expressed
under both the merged-matrix and the per-study-effect-size views
is more credible than any gene found under one view alone.

\subsection{Matrix dissection: HarmonizR}

Vo{\ss} et al.~\cite{voss2022} introduced HarmonizR, a framework
that applies batch correction to incomplete matrices \emph{without}
imputation. The key idea is \emph{matrix dissection}: the combined
data matrix is split into sub-matrices that share a common pattern
of non-missing values (i.e.\ every feature in a sub-matrix is
observed in every sample of that sub-matrix). ComBat or limma's
\texttt{removeBatchEffect()} is then applied independently to each
sub-matrix, and the corrected sub-matrices are reassembled into the
full matrix. Features that are observed in only a single batch
cannot be corrected and are returned unchanged.

HarmonizR was developed for proteomics, where block-structured
missingness is endemic: different LC-MS/MS acquisition modes (DDA,
DIA, SWATH) yield overlapping but non-identical protein inventories
across batches, exactly paralleling the gene-coverage problem in
our microarray collection. Evaluated on four proteomic datasets
with up to 23 batches, HarmonizR outperformed the impute-then-ComBat
strategy in terms of significant-protein detection while
preserving within-batch statistical properties~\cite{voss2022}.
A follow-up paper~\cite{voss2025} added blocking and singular-feature
handling, further improving runtime and data preservation.

The method is available as a Bioconductor R package
(\texttt{HarmonizR}). For the placenta collection it is directly
applicable: the merged expression matrix can be fed to HarmonizR
without prior imputation, and the framework will automatically
partition genes by their dataset-coverage pattern and apply ComBat
to each partition. This approach avoids the chicken-and-egg problem
of needing batch correction before imputation but imputation before
batch correction.

\subsection{Batch-sensitized imputation}

Goh, Hui, and Wong~\cite{goh2023} formalised the interaction
between missing value imputation (MVI) and batch effect correction
(BEC). They compared three imputation strategies:
\begin{itemize}
  \item \textbf{M1 (global):} impute using information from all
        samples regardless of batch;
  \item \textbf{M2 (self-batch / batch-sensitized):} split by batch,
        impute within each batch separately;
  \item \textbf{M3 (cross-batch):} impute using only samples from
        the \emph{opposite} batch.
\end{itemize}
On simulated and real proteomics/genomics data, M2 consistently
outperformed M1 and M3: batch-sensitized imputation preserved
within-batch variance structure and produced lower RMSE after
downstream ComBat correction. Crucially, both M1 and M3 increased
sample variance following ComBat, whereas M2 did not.

The authors also introduced the concept of \emph{batch effect
associated missing values} (BEAMs) --- missingness that is itself
correlated with the batch factor. When BEAMs are present (as in
our microarray collection, where gene coverage differs by platform),
global imputation (M1) is particularly problematic because the
imputed values carry no batch-specific information, making subsequent
BEC less effective.

For the placenta pipeline, this motivates a modification of the
current imputation step: rather than applying softImpute to the
full merged matrix (a global M1 strategy), imputation could be
performed per-dataset (M2) before merging. However, per-dataset
imputation is only possible for genes measured in at least some
samples within that dataset; genes entirely absent from a dataset
cannot be imputed within it, so the coverage problem remains and
must still be addressed by cross-dataset methods such as HarmonizR
or the current softImpute approach.

No standalone R package was released with the Goh et al.\ paper;
the M2 strategy is a simple procedural change (split by batch,
impute, recombine) applicable to any imputation method.

\subsection{Imputation as a complement, not a replacement}

The main manuscript advocates matrix-completion imputation
(\textit{softImpute}, $k$-NN) followed by ComBat. This combination
addresses the \emph{coverage} problem (different platforms target
different gene catalogues) but not the \emph{design imbalance}
problem (different studies cover different trimesters). The two
are independent: imputing missing genes does not give ComBat any
new information about the trimester--batch confound. Any of the
methods in this supplement can be applied to the imputed matrix
exactly as they would be applied to the raw merged matrix.

%% =====================================================================
\section{Recommendations for the placenta pipeline}

We suggest the following layered strategy, ordered from cheapest to
most invasive, so that each layer can be added without rebuilding
the previous one:

\begin{enumerate}
  \item \textbf{Document the imbalance.} Tabulate, for every
        contrast of interest, which studies contribute samples from
        which trimester. Any contrast for which one trimester comes
        from disjoint studies is intrinsically risky and should be
        flagged in the results.

  \item \textbf{Replace ``ComBat $\to$ limma'' with
        ``limma with batch in the design.''} This is the
        Nygaard/Towfic recommendation~\cite{nygaard2016,towfic2017}.
        It is one line of R and removes the largest source of bias
        for free, at the cost of some power when the design happens
        to be balanced.

  \item \textbf{Run RUV with negative-control genes
        in parallel.} Use the imputed, merged expression matrix as
        input. Restrict to housekeeping genes intersected with
        empirically invariant genes from a balanced reference
        study. Compare the resulting DEG list against the
        ComBat-in-design list and report the intersection.

  \item \textbf{Run an effect-size meta-analysis as a sanity
        check.} Compute per-study log fold changes for every
        contrast in which a study contains both endpoints, then
        combine with a random-effects model (\texttt{metafor}).
        Report only genes that are significant under both the
        joint-matrix and the meta-analytic views.

  \item \textbf{If a single dataset spans the broadest range of
        trimesters, try reference-batch ComBat (or ComBat-ref) with
        that dataset as the anchor.} Use this as a third comparison
        point, not as the primary correction.

  \item \textbf{Avoid SVA on this collection unless the residual
        structure has been inspected first.} The risk that SVA
        absorbs trimester into a surrogate variable is real and
        hard to detect without a balanced ground-truth subset.

  \item \textbf{Consider HarmonizR as an alternative to
        imputation + ComBat.} HarmonizR~\cite{voss2022} applies
        ComBat to sub-matrices sharing common coverage patterns,
        bypassing imputation entirely. This avoids the
        imputation--batch-correction confound documented by
        Goh et al.~\cite{goh2023} and is directly applicable to
        the block-structured missingness in the placenta collection.

  \item \textbf{If imputation is used, prefer batch-sensitized
        imputation (M2).} Following Goh et al.~\cite{goh2023},
        impute within each dataset separately before merging, rather
        than imputing on the combined matrix. This preserves
        within-batch variance structure and improves subsequent
        batch correction.
\end{enumerate}

The above strategy preserves the existing imputation step --- which
addresses a different problem (cross-platform gene coverage) ---
and supplements ComBat with methods that do not require a balanced
design.

%% =====================================================================
\begin{thebibliography}{99}

\bibitem{johnson2007}
Johnson WE, Li C, Rabinovic A.
Adjusting batch effects in microarray expression data using
empirical Bayes methods.
\textit{Biostatistics} 2007;8(1):118--127.
\doi{10.1093/biostatistics/kxj037}

\bibitem{nygaard2016}
Nygaard V, R{\o}dland EA, Hovig E.
Methods that remove batch effects while retaining group differences
may lead to exaggerated confidence in downstream analyses.
\textit{Biostatistics} 2016;17(1):29--39.
\doi{10.1093/biostatistics/kxv027}

\bibitem{towfic2017}
Towfic F, Kusko R, Zeskind B.
Letter to the Editor response: Nygaard et al.
\textit{Biostatistics} 2017;18(2):195--197.
\doi{10.1093/biostatistics/kxw031}

\bibitem{leek2007}
Leek JT, Storey JD.
Capturing heterogeneity in gene expression studies by surrogate
variable analysis.
\textit{PLoS Genetics} 2007;3(9):e161.
\doi{10.1371/journal.pgen.0030161}

\bibitem{gagnonbartsch2012}
Gagnon-Bartsch JA, Speed TP.
Using control genes to correct for unwanted variation in
microarray data.
\textit{Biostatistics} 2012;13(3):539--552.
\doi{10.1093/biostatistics/kxr034}

\bibitem{risso2014}
Risso D, Ngai J, Speed TP, Dudoit S.
Normalization of RNA-seq data using factor analysis of control
genes or samples.
\textit{Nature Biotechnology} 2014;32(9):896--902.
\doi{10.1038/nbt.2931}

\bibitem{walker2008}
Walker WL, Liao IH, Gilbert DL, et al.
Empirical Bayes accommodation of batch effects in microarray data
using identical replicate reference samples.
\textit{BMC Genomics} 2008;9:494.
\doi{10.1186/1471-2164-9-494}

\bibitem{varma2020}
Varma S.
Blind estimation and correction of microarray batch effect.
\textit{PLoS ONE} 2020;15(4):e0231446.
\doi{10.1371/journal.pone.0231446}

\bibitem{zhang2020}
Zhang Y, Parmigiani G, Johnson WE.
ComBat-seq: batch effect adjustment for RNA-seq count data.
\textit{NAR Genomics and Bioinformatics} 2020;2(3):lqaa078.
\doi{10.1093/nargab/lqaa078}

\bibitem{zhang2025}
Zhang X.
Highly effective batch effect correction method for RNA-seq count
data (ComBat-ref).
\textit{Computational and Structural Biotechnology Journal}
2025;27:58--68.
\doi{10.1016/j.csbj.2024.12.010}

\bibitem{haghverdi2018}
Haghverdi L, Lun ATL, Morgan MD, Marioni JC.
Batch effects in single-cell RNA-sequencing data are corrected by
matching mutual nearest neighbors.
\textit{Nature Biotechnology} 2018;36(5):421--427.
\doi{10.1038/nbt.4091}

\bibitem{korsunsky2019}
Korsunsky I, Millard N, Fan J, et al.
Fast, sensitive and accurate integration of single-cell data with
Harmony.
\textit{Nature Methods} 2019;16(12):1289--1296.
\doi{10.1038/s41592-019-0619-0}

\bibitem{turnbull2012}
Turnbull AK, Kitchen RR, Larionov AA, et al.
Direct integration of intensity-level data from Affymetrix and
Illumina microarrays improves statistical power for robust
reanalysis.
\textit{BMC Medical Genomics} 2012;5:35.
\doi{10.1186/1755-8794-5-35}

\bibitem{macedo2013}
Almeida-de-Macedo MM, Ransom N, Feng Y, Hurst J, Wurtele ES.
Comprehensive analysis of correlation coefficients estimated from
pooling heterogeneous microarray data.
\textit{BMC Bioinformatics} 2013;14:214.
\doi{10.1186/1471-2105-14-214}

\bibitem{shi2011}
Shi F, Abraham G, Leckie C, Haviv I, Kowalczyk A.
Meta-analysis of gene expression microarrays with missing
replicates.
\textit{BMC Bioinformatics} 2011;12:84.
\doi{10.1186/1471-2105-12-84}

\bibitem{voss2022}
Vo{\ss} H, Schlumbohm S, Barwikowski P, et al.
HarmonizR enables data harmonization across independent proteomic
datasets with appropriate handling of missing values.
\textit{Nature Communications} 2022;13:3523.
\doi{10.1038/s41467-022-31007-x}

\bibitem{voss2025}
Vo{\ss} H, Schlumbohm S, Barwikowski P, et al.
HarmonizR: blocking and singular feature data adjustment improve
runtime efficiency and data preservation.
\textit{BMC Bioinformatics} 2025;26:73.
\doi{10.1186/s12859-025-06073-9}

\bibitem{goh2023}
Goh WWB, Hui HWH, Peng H, Wong L.
How missing value imputation is confounded with batch effects and
what you can do about it.
\textit{Drug Discovery Today} 2023;28(9):103661.
\doi{10.1016/j.drudis.2023.103661}

%% --- Source datasets ---

\bibitem{GSE100051}
Soncin F, Khater M, To C, et al.
Comparative analysis of mouse and human placentae across gestation
reveals species-specific regulators of placental development.
\textit{Development} 2018;145(2):dev156273.
\doi{10.1242/dev.156273}

\bibitem{GSE122214}
Zhao L, Sun LF.
The placental transcriptome of the first-trimester placenta is
affected by in vitro fertilization and embryo transfer.
\textit{Reproductive Biology and Endocrinology} 2019;17:54.
\doi{10.1186/s12958-019-0494-7}

\bibitem{GSE22490}
Rull K, Nagirnaja L, Ulander VM, et al.
Increased placental expression and maternal serum levels of
apoptosis-inducing TRAIL in recurrent miscarriage.
\textit{Placenta} 2013;34(2):141--148.
\doi{10.1016/j.placenta.2012.11.032}

\bibitem{GSE28551}
Sitras V, Paulssen RH, Gr{\o}naas H, et al.
Differences in gene expression between first and third trimester
human placenta: a microarray study.
\textit{PLoS ONE} 2012;7(3):e33294.
\doi{10.1371/journal.pone.0033294}

\bibitem{GSE35574}
Guo L, Tsai SQ, Hardison NE, et al.
Differentially expressed microRNAs and affected biological pathways
revealed by modulated modularity clustering (MMC) analysis of human
preeclamptic and IUGR placentas.
\textit{Placenta} 2013;34(7):599--605.
\doi{10.1016/j.placenta.2013.04.007}

\bibitem{GSE37653}
Khan MA, Manna S, Engstr{\"o}m K, et al.
Expressional regulation of genes linked to immunity \&
programmed development in human early placental villi.
\textit{Indian Journal of Medical Research} 2014;139(1):125--140.
PMID:~24604048.

\bibitem{GSE37901}
Uusk{\"u}la L, R{\~o}{\~o}t A, Reimand J, et al.
Mid-gestational gene expression profile in placenta and link to
pregnancy complications.
\textit{PLoS ONE} 2012;7(11):e49248.
\doi{10.1371/journal.pone.0049248}

% GSE55439: no associated publication (GEO citation missing).
% Submitters: Wolfe LM, Thagarajan RD, Boscolo F, et al.

\bibitem{GSE6573}
Herse F, Dechend R, Harsem NK, et al.
Dysregulation of the circulating and tissue-based
renin-angiotensin system in preeclampsia.
\textit{Hypertension} 2007;49(3):604--611.
\doi{10.1161/01.HYP.0000257797.49289.71}

\bibitem{GSE70102}
Bianco K, Gormley M, Engel N, et al.
Placental transcriptomes in the common aneuploidies reveal
critical regions on the trisomic chromosomes and genome-wide
effects.
\textit{Prenatal Diagnosis} 2016;36(9):812--822.
\doi{10.1002/pd.4862}

\bibitem{GSE73374}
Martin E, Smeester L, Bommarito PA, et al.
Epigenetics and preeclampsia: defining functional epimutations in
the preeclamptic placenta related to the TGF-$\beta$ pathway.
\textit{PLoS ONE} 2015;10(10):e0141294.
\doi{10.1371/journal.pone.0141294}

\bibitem{GSE73685}
Bukowski R, Sadovsky Y, Goodarzi H, et al.
Onset of human preterm and term birth is related to unique
inflammatory transcriptome profiles at the maternal fetal
interface.
\textit{PeerJ} 2017;5:e3685.
\doi{10.7717/peerj.3685}

\bibitem{GSE93520}
Zhang B, Zheng Z, Wu X, et al.
Villi-specific gene expression reveals novel prognostic
biomarkers in multiple human cancers.
\textit{Journal of Cancer} 2017;8(16):3193--3200.
\doi{10.7150/jca.19787}

\bibitem{GSE9984}
Mikheev AM, Nabekura T, Kaddoumi A, et al.
Profiling gene expression in human placentae of different
gestational ages: an OPRU Network and UW SCOR Study.
\textit{Reproductive Sciences} 2008;15(9):866--877.
\doi{10.1177/1933719108322425}

\end{thebibliography}

\end{document}
